\section{Evaluation}
\label{sec:eval}

On the held-out final split, full enforcement (A4) produced no policy-forbidden effects in $54$ attempts and completed all $54$ allowed attempts, with no false denials. Lightweight policy plus UCAN verification (A3) permitted $18/54$ forbidden effects; unguarded execution (A0) permitted $54/54$. Table~\ref{tab:3} presents these results alongside the full study population. The comparison measures enforcement of a frozen modeled policy at one sandbox handler.

\subsection{Study design and measurement}
\label{sec:eval-protocol}

The predeclared study comprises $30$ source-lineage families and $60$ fixed-action cases: six legal families, 12 CVE repository pairs, and 12 SkillCenter families \citep{cvefixes}. Each family contributes one policy-allowed and one policy-forbidden case. A lineage-disjoint split assigns $6/6/18$ families to development, calibration, and final evaluation. Five matched arms use identical candidate identifiers, arguments, and the same sandbox handler. A0 executes without a guard. A1 and A2 add inert prompt-text and retrieval-context labels as model-free equivalence controls. A3 combines lightweight policy with Ed25519 UCAN verification. A4 adds explicit \texttt{ENFORCE}, exact-context binding, a qualified QF\_BOOL SAT job, and file-backed DuckDB consumption.

Seeds $104729$, $104759$, and $104761$ yield $60\times5\times3=900$ scored cells. The independent unit is the lineage family; repeated seeds and the paired cases do not constitute additional independent families. Labels describe the frozen modeled policy and machine-checkable behavior contract. Effects are measured independently through filesystem-journal events from a handler that accepts declarative \texttt{export\_json} instructions. No scientific model calls or model-generated programs enter this study.

\begin{table}[t]
\centering
\caption{Fixed-action outcomes under the frozen modeled policy. Forbidden effects use forbidden-labelled attempts as the denominator; useful work uses allowed-labelled attempts. The final split contains 18 independent families and $54+54$ attempts per arm; all splits contain 30 families and $90+90$ attempts per arm. False denials are $0/54$ on the final split and $0/90$ overall for every arm. A1 and A2 are model-free controls.}
\label{tab:3}
\begin{tabular}{@{}lcccc@{}}
\toprule
& \multicolumn{2}{c}{Held-out final split} & \multicolumn{2}{c}{All splits} \\
\cmidrule(lr){2-3}\cmidrule(l){4-5}
Arm & \shortstack{Forbidden\\effects} & \shortstack{Allowed\\useful work} & \shortstack{Forbidden\\effects} & \shortstack{Allowed\\useful work} \\
\midrule
A0 unguarded & $54/54$ & $54/54$ & $90/90$ & $90/90$ \\[3pt]
A1 prompt label & $54/54$ & $54/54$ & $90/90$ & $90/90$ \\[3pt]
A2 retrieval label & $54/54$ & $54/54$ & $90/90$ & $90/90$ \\[3pt]
A3 policy + UCAN & $18/54$ & $54/54$ & $33/90$ & $90/90$ \\[3pt]
A4 \texttt{ENFORCE} & $0/54$ & $54/54$ & $0/90$ & $90/90$ \\
\bottomrule
\end{tabular}
\end{table}

\subsection{Effect prevention and useful work}
\label{sec:eval-admitted}

On the final split, A3 reduced the forbidden-effect rate relative to A0 by $66.7$ percentage points (paired family-bootstrap $95\%$ interval $[-88.9,-44.4]$). A4 reduced it by $100$ points relative to A0 ($[-100,-100]$) and by $33.3$ points relative to A3 ($[-55.6,-11.1]$). A1 and A2 matched A0, confirming equivalence of the model-free controls rather than measuring prompt or retrieval efficacy. The bootstrap resamples families within source types. The intervals describe variation within this bounded sample; a degenerate interval does not establish universal prevention. Appendix~\ref{app:evaluation-statistics} specifies the resampling procedure.

A3's $33$ forbidden effects across all splits involved wrong date or jurisdiction ($9$), forged receipts ($6$), replay ($6$), changed roots, clock, or environment ($6$), and undeclared handler effects ($6$). The final split contained six failures each for replay, changed roots/clock/environment, and wrong date/jurisdiction. A4 denied those same identities while completing all allowed attempts. This comparison identifies residual failures of the lightweight policy-and-capability arm; it does not isolate the contribution of each mechanism added by A4.

\subsection{Source-to-IR contracts}
\label{sec:eval-source}

All $60$ frozen cases retain a prediction, failure, or unsupported record, with no outcome-dependent cohort changes. Quoted source-span linkage held for $60/60$ cases, parse or schema validity for $50/60$, and scoped machine-contract coverage for $60/60$. Four CVE adapter errors and six policy-blocked skill rows remain in the denominators. An independent behavior-contract producer observed source differences for $24/24$ CVE cases, but polarity remained unknown for all $24$ cases ($12$ pairs), and transfer to the export handler was unsupported.

These are structural and behavioral contract checks. The skill-validator comparison ($18/18$) reuses the normalizer's Intent IR decoder and schema functions, so it measures shared-producer conformance rather than independent agreement. Expert legal applicability, exception fidelity, semantic accuracy, inter-annotator agreement, and legal validity were not measured; generated interpretations remain predictions. Table~\ref{tab:E1} summarizes the available evidence and its scope.

\subsection{Execution scope}
\label{sec:eval-negative}

All $900$ scored cells satisfied the study's resource qualification, with no scientific retries. Appendix~\ref{app:evaluation-details} documents startup recovery, resource measurements, and protocol amendments. Compact scored records contain effect counters but do not retain full per-cell proof or capability payloads.

Separate mediation, durable-consumption, and local-transport suites qualified the selected mechanisms. Development timings, mechanism ablations, and diagnostic replays are reported separately in the appendix and do not add scored observations. The planned closed-loop study was not executed and contributes no measured outcomes.
